\documentclass[10pt,a4paper]{article}
\usepackage[a4paper,left=1.6cm,right=1.6cm,top=2.0cm,bottom=1.6cm]{geometry}
\usepackage[T1]{fontenc}
\usepackage[scaled=0.94]{helvet}
\renewcommand{\familydefault}{\sfdefault}
\usepackage{amsmath,amssymb}
\usepackage[table]{xcolor}
\usepackage{booktabs,longtable,array,tabularx}
\usepackage{fancyhdr,titlesec,enumitem,fancyvrb}
\usepackage{upquote}
\usepackage[most]{tcolorbox}
\definecolor{ink}{HTML}{1B2A41}
\definecolor{accent}{HTML}{1F5F8B}
\definecolor{okgreen}{HTML}{2E7D4F}
\usepackage[colorlinks=true,linkcolor=accent,urlcolor=accent]{hyperref}
\usepackage[expansion=false,protrusion=false]{microtype}
\DisableLigatures[-]{encoding = T1, family = tt*}

\input{variant}

\titleformat{\section}{\Large\bfseries\color{ink}}{\thesection.}{0.6em}{}
\titleformat{\subsection}{\large\bfseries\color{accent}}{\thesubsection}{0.6em}{}
\titlespacing*{\section}{0pt}{2.2ex plus .5ex}{0.8ex}
\titlespacing*{\subsection}{0pt}{1.6ex plus .4ex}{0.5ex}
\pagestyle{fancy}
\fancyhf{}
\fancyhead[L]{\small\color{accent}xphd: how to run}
\fancyhead[R]{\small\color{gray}xphd \xphdversion}
\fancyfoot[R]{\small\color{gray}\thepage}
\renewcommand{\headrulewidth}{0.4pt}
\setlength{\headheight}{15pt}
\setlength{\parskip}{0.55ex}
\setlength{\parindent}{0pt}
\setlength{\LTpre}{4pt}\setlength{\LTpost}{6pt}

\newcommand{\xphd}{\texttt{xphd}}
\newcommand{\hd}[1]{\textcolor{white}{\normalfont\bfseries #1}}
\newcommand{\hdrow}{\rowcolor{accent}}
\newcommand{\bQ}{\mathbf{Q}}
\newcommand{\bq}{\mathbf{q}}
\newtcblisting{cmd}{listing only,breakable,colback=black!4,colframe=black!18,
  boxrule=0.4pt,arc=1pt,left=4pt,right=4pt,top=2pt,bottom=2pt,
  listing options={upquote=true,basicstyle=\ttfamily\footnotesize,breaklines=true,columns=fullflexible,keepspaces=true,breakindent=1.5em}}
\newtcolorbox{note}[1][]{colback=okgreen!7,colframe=okgreen!70!black,boxrule=0.5pt,arc=2pt,
  left=6pt,right=6pt,top=4pt,bottom=4pt,breakable,#1}
\newcommand{\step}[4]{%
  {\small\noindent\begin{tabularx}{\linewidth}{@{}>{\bfseries\leavevmode\color{accent}}p{2.0cm}X@{}}
  What & #1\\ Needs & #2\\ Writes & #3\\ Check & #4\end{tabularx}\par}\vspace{1.0ex}}

\begin{document}
\thispagestyle{fancy}
{\color{ink}\fontsize{26}{30}\selectfont\bfseries xphd: how to run\par}
\vspace{0.8ex}
{\raggedright\color{accent}\large Exciton--phonon matrix elements, linewidths, symmetry, transport and photoemission from Quantum~ESPRESSO, LetzElPhC and yambo\par}
\vspace{1.0ex}
{Rishabh Saraswat, Mantasha Ansari, Rekha Verma and Sitangshu Bhattacharya\par}
{\small Indian Institute of Information Technology, Allahabad \quad\textperiodcentered\quad Version \xphdversion, MIT License\par}
{\small Cite: R.~Saraswat, M.~Ansari, R.~Verma and S.~Bhattacharya, arXiv:2610.12996 (2026); xphd, \xphdsoftware\par}
\vspace{1.5ex}

\begin{note}
This is the only document. It lists what the upstream codes must produce, every command in the order it is run, what each
needs and writes, and how to tell that its output is right. Every option of every command is printed by
\texttt{xphd <command> --help} and collected in the reference at the end. Every run writes a report next to its output, a PDF
on Windows and a text file on Linux (Sec.~\ref{sec:report}).
\end{note}

{\setlength{\parskip}{0pt}\small\tableofcontents}

% =====================================================================
\section{What xphd needs and what it gives}
\begin{tabularx}{\linewidth}{>{\bfseries}p{3.2cm}X}
Needs & a yambo \texttt{SAVE} (from \texttt{p2y} on the nscf); the BSE at every irreducible $\bQ$ with full diagonalisation
(\texttt{ndb.BS\_diago\_Q*}); LetzElPhC's \texttt{ndb.elph} built on \emph{the same} \texttt{SAVE}, over the BSE band window;
\texttt{matdyn.x} frequencies (and eigenvectors) on a fine $q$ mesh; optionally the GW \texttt{ndb.QP} and \texttt{ndb.dipoles}.\\
Gives & exciton--phonon matrix elements $G(\bQ,\bq)$ on the full zone, one archive per $\bQ$; broadening-free linewidths,
self-energies and absorption lineshapes; the mirror ($\sigma_h$) parity and irreducible representations of excitons, bands
and phonons, and the selection rules they imply; exciton dipoles and chirality; real-time Boltzmann transport of the exciton
populations; exciton photoemission maps (thermal or time-resolved).\\
\end{tabularx}

% =====================================================================
\section{Install}
\begin{cmd}
cd xphd
pip install -e ".[plot,yambo]"      # numpy, matplotlib, netCDF4
pip install <path to yambopy>        # from source; needed by every command that reads yambo or LetzElPhC
pip install -e ".[hpc]"              # clusters: mpi4py (built against the cluster MPI), scipy, threadpoolctl
xphd --help                          # every command
xphd <command> --help                # its options
python -m pytest                     # synthetic data; no yambo databases needed
\end{cmd}
Only \texttt{numpy} is required to import the package; the other libraries are imported by the commands that need them.
\texttt{scipy} is needed by \texttt{bte}. Paths in angle brackets are placeholders.

% =====================================================================
\section{The workflow at a glance}
{\small\renewcommand{\arraystretch}{1.25}
\begin{tabularx}{\linewidth}{>{\bfseries\raggedright\arraybackslash}p{2.3cm}>{\raggedright\arraybackslash}p{5.8cm}>{\raggedright\arraybackslash}X}
\hdrow \hd{Stage} & \hd{Code} & \hd{Output}\\
A. Ground state & \texttt{pw.x} scf, nscf; \texttt{p2y} & \texttt{SAVE}\\
B. Phonons & \texttt{ph.x}; \texttt{q2r.x}; \texttt{matdyn.x} & dyn, dvscf; force constants; fine-mesh frequencies\\
C. Excitons & yambo BSE (GW optional) & \texttt{ndb.BS\_diago\_Q1} \ldots\ at every irreducible $\bQ$\\
C$'$. El--ph & LetzElPhC & \texttt{ndb.elph} (same \texttt{SAVE} as C)\\
D. Archives & \texttt{xphd dmats}, \texttt{generate}, \texttt{check-archive}, \texttt{verify-archives} & \texttt{Dmats.npy}, \texttt{GI\_ExcPh\_Q0001.npz} \ldots\\
E. Fine phonons & \texttt{xphd matdyn-split}, \texttt{matdyn-merge}, \texttt{modes-check}, \texttt{matdyn}, \texttt{mode-labels} & \texttt{hw\_fine.npy}, \texttt{mode\_labels.npy}\\
F. Widths & \texttt{xphd linewidth}, \texttt{sweep}, \texttt{unfold}, \texttt{selfenergy}, \texttt{lineshape} & widths, zone maps, $\Sigma(\omega)$, $\varepsilon_2$\\
G. Symmetry & \texttt{xphd parity}, \texttt{band-parity}, \texttt{irreps}, \texttt{selection-rule} & parity maps, labels, allowed/forbidden split\\
H. Optics & \texttt{xphd dipoles}, \texttt{chirality} & \texttt{exc\_dipoles.npy}, helicity\\
I. Dynamics & \texttt{xphd bte}, \texttt{arpes} & snapshots; photoemission maps\\
\end{tabularx}}

D needs A--C$'$; E needs B; F needs D and E; G needs A, C and D; I needs D and E (and G for \texttt{bte --window}).

% =====================================================================
\section{Upstream settings}
xphd trusts its inputs; these must hold.
\subsection{Quantum ESPRESSO}
\begin{itemize}[nosep,leftmargin=1.4em]
\item scf, phonons and nscf in separate directories (the nscf overwrites \texttt{<prefix>.save}).
\item q-grid = k-grid ($N\times N\times1$), so every $\mathbf{k}+\bq$ lies on the mesh. Set \texttt{fildvscf} in \texttt{ph.x}.
\item A \texttt{ph.x} restart with \texttt{recover=.true.} trusts the records in \texttt{\_ph0/<prefix>.phsave}: remove the
\texttt{dynmat.<iq>.*.xml} of any $\bq$ to be recomputed first, or its dvscf is left empty. Check that every dvscf has the same size.
\item Spin--orbit: \texttt{noncolin=.true.}, \texttt{lspinorb=.true.}, fully relativistic pseudopotentials; \texttt{ph.x} then writes spinor dvscf.
\end{itemize}
\subsection{q2r.x and matdyn.x}
\begin{cmd}
 q2r.in                              matdyn.in.tmpl  (namelist only; the q-list is appended per chunk)
 &input                              &input
   fildyn  = '<prefix>.dyn'            asr = 'crystal', flfrc = '<prefix>.fc'
   zasr    = 'crystal'                 flfrq = 'FLFRQ', flvec = 'FLVEC'
   flfrc   = '<prefix>.fc'             loto_2d = .true., q_in_cryst_coord = .true.
   loto_2d = .true.                    q_in_band_form = .false., nosym = .true., dos = .false.
 /                                   /
\end{cmd}
\texttt{loto\_2d=.true.} in both (the 3D non-analytic term adds a spurious LO splitting at $\Gamma$). The fine mesh is
$N_f = N\times r$ with $r$ the refinement used later (e.g.\ $18\times20=360$).
\subsection{LetzElPhC}
\begin{cmd}
nqpool      = 4               # nqpool x nkpool = the MPI ranks
nkpool      = 12
start_bnd   = <first BSE band>
end_bnd     = <last BSE band>
save_dir    = <SAVE>          # the SAVE the BSE reads
ph_save_dir = ph.save         # the default is ph_save; misspelt keys are ignored silently
kernel      = dfpt
convention  = yambo
\end{cmd}
\texttt{ph.save} gathers \texttt{dyn0..N}, \texttt{dvscf1..N}, \texttt{patterns.N.xml}, \texttt{tensors.xml},
\texttt{data-file-schema.xml} and the pseudopotentials (LetzElPhC's \texttt{lelphc -pp}). Run \texttt{mpirun -np 48 lelphc -F lelphc.in}.
\textbf{Same SAVE:} LetzElPhC and the BSE must read the same \texttt{SAVE} (one wavefunction gauge); redo both if the nscf is redone.
The precision of LetzElPhC's build must match that of the \texttt{SAVE} (single or double).
\subsection{yambo}
\begin{itemize}[nosep,leftmargin=1.4em]
\item 2D: \texttt{CUTGeo= "slab z"} with the RIM. BSE: \texttt{BSEBands} is the band window (it fixes \texttt{--nv}),
\texttt{BSEQptR} covers every irreducible $\bQ$, \texttt{BSSmod= "d"} (full diagonalisation) and \texttt{WRbsWF} (eigenvectors written).
\item Check that the energies in \texttt{ndb.BS\_diago\_Q1} are not all zero and that the report's \texttt{[K+QP]} gap equals the
Kohn--Sham gap plus the GW correction.
\end{itemize}

% =====================================================================
\section{Commands, in order}
\subsection{Archives}
\begin{cmd}
xphd ibz --savepath <SAVE>
xphd dmats --bse-dir <BSE_DIR> --savepath <SAVE> --ndb-elph <ELPH> --nexc <n>
xphd generate 0 <ARCH> --nexc <n> --savepath <SAVE> --ndb-elph <ELPH> --bse-dir <BSE_DIR> --dmats Dmats.npy
xphd check-archive <ARCH>/GI_ExcPh_Q0001.npz --detail
mpirun -np <R> xphd generate all <ARCH> --nexc <n> --savepath <SAVE> --ndb-elph <ELPH> \
       --bse-dir <BSE_DIR> --dmats Dmats.npy
xphd verify-archives --dir <ARCH> --n <N*N>
\end{cmd}
\step{\texttt{ibz} lists the irreducible $\bQ$. \texttt{dmats} builds the single-particle rotation matrices over the BSE band window.
\texttt{generate} computes $G(\bQ,\bq)$ for every $\bq$, branch and pair of states, rotating excitons from the irreducible BSE;
\texttt{0} is $\Gamma$, \texttt{all} the whole zone (shared between MPI ranks, resumable).}
{\texttt{SAVE}, the BSE databases, \texttt{ndb.elph}.}
{\texttt{Dmats.npy}; \texttt{GI\_ExcPh\_Q<iQ+1>.npz} per $\bQ$.}
{\texttt{check-archive} prints READY, with time reversal and $C_3$ both near $10^{-3}$; one far above the other is a rotation fault
(or a split doublet: \texttt{--detail}, \texttt{--deg-tol}). \texttt{verify-archives}: all $N^2$ present and complete.}

\subsection{Phonons on the fine mesh}
\begin{cmd}
q2r.x < q2r.in
xphd matdyn-split --n <Nf> --chunks 48 --prefix qlist
for c in $(seq -f %03g 0 47); do                       # matdyn.x is serial: one process per chunk
  sed -e "s/FLVEC/modes_$c.modes/" -e "s/FLFRQ/freq_$c.freq/" matdyn.in.tmpl > in_$c
  cat qlist_$c.txt >> in_$c; matdyn.x < in_$c > out_$c &
done; wait
xphd matdyn-merge --chunks 48 -o matdyn.modes --freq matdyn.freq
xphd modes-check --modes matdyn.modes --freq matdyn.freq --n <Nf>
xphd matdyn matdyn.freq --fine <Nf> --nbnd <3 x atoms> --npz <ARCH>/GI_ExcPh_Q0001.npz -o hw_fine.npy
xphd mode-labels --modes matdyn.modes --fine <Nf> --mesh <N> --masses <m_1> <m_2> \
                 --archive <ARCH>/GI_ExcPh_Q0001.npz --hw-fine hw_fine.npy -o mode_labels.npy
\end{cmd}
\step{The fine-mesh frequencies in eV, row-major; \texttt{mode-labels} gives each branch its in-/out-of-plane character on the archive mesh
(planar two-atom layers; for layers where $\sigma_h$ exchanges atoms take the parities from the el--ph database).}
{The force constants; one archive for its $\bq$ ordering.}
{\texttt{matdyn.modes}, \texttt{matdyn.freq}, \texttt{hw\_fine.npy}, \texttt{mode\_labels.npy}.}
{\texttt{modes-check} passes; \texttt{matdyn} reproduces the archive's frequencies where the meshes coincide ($\mu$eV).
Negative frequencies only next to $\Gamma$ in the flexural branch.}

\subsection{Linewidths, self-energy and lineshape}
\begin{cmd}
xphd linewidth <ARCH>/GI_ExcPh_Q0001.npz --hw-fine hw_fine.npy --T 10 77 300 \
      --acoustic-model q2 --n-shell 1 --per-mode -o G.npz
xphd sweep '<ARCH>/GI_ExcPh_Q*.npz' --hw-fine hw_fine.npy --T 77 --acoustic-model q2 --n-shell 1 \
      --workers 8 --unfold <SAVE> -o lw_77K.npz
xphd selfenergy <ARCH>/GI_ExcPh_Q0001.npz --hw-fine hw_fine.npy --state 1 --T 300 -o se.npz
xphd lineshape <ARCH>/GI_ExcPh_Q0001.npz --T 10 77 300 --hw-fine hw_fine.npy \
      --acoustic-model shell --acoustic-scaling omega --parity parity.npz \
      --bsdb <BSE_DIR>/ndb.BS_diago_Q1 --n-omega 4801 -o lineshape.npz
\end{cmd}
\step{\texttt{linewidth}: golden-rule width of every state at one $\bQ$, the energy $\delta$ integrated on fine-mesh triangles
(exciton energies refined by local quadratics), emission/absorption and per branch. \texttt{sweep}: the same at every $\bQ$, unfolded
onto the zone. \texttt{selfenergy}: $\Sigma(\omega)$ of one state with its Kramers--Kronig partner. \texttt{lineshape}: the full
$\Sigma$ matrix and the absorption line; final states carry the self-consistent width (\texttt{--final-width sc}, default).}
{Archives and \texttt{hw\_fine.npy}; for \texttt{lineshape} the dipoles (\texttt{--bsdb} or \texttt{--dipoles}) and \texttt{parity.npz}.}
{\texttt{G.npz}, \texttt{lw\_77K.npz}, \texttt{se.npz}, \texttt{lineshape.npz}.}
{The acoustic channel inside the cell around $q=0$ is zero by the sum rule unless a model is chosen: \texttt{none} drops it, which
underestimates the width at a band minimum. \texttt{lineshape} prints the symmetry check (opposite parities $\approx0$, degenerate blocks
proportional to the identity), the final-width gain (well below 1) and warns if the line is narrower than half its on-shell width; the widths
at low $T$ need a frequency step below the width (\texttt{--n-omega}).}

\subsection{Symmetry}
\begin{cmd}
xphd parity --save <SAVE> --bse-dir <BSE_DIR> --dmats Dmats.npy --nv <nv> --nstates <n> \
            --deg-tol 0.005 -o parity.npz
xphd band-parity --save <SAVE> --bse-dir <BSE_DIR> --dmats Dmats.npy --nv <nv> --path G M K G -o band_parity.npz
xphd irreps --save <SAVE> --bse-dir <BSE_DIR> --dmats Dmats.npy --nv <nv> --points G K M \
            --nstates <n> -o irreps.npz
xphd selection-rule --archive <ARCH>/GI_ExcPh_Q0001.npz --parity parity.npz --labels mode_labels.npy \
                    --state 1 --by-final
\end{cmd}
\step{$\chi(\sigma_h)$ of every exciton at every $\bQ$ (unfolded with yambo's own map); band parities; irrep labels with the closure
of each multiplet; the coupling split into symmetry-allowed and -forbidden.}
{\texttt{SAVE}, the BSE at all $\bQ$, \texttt{Dmats.npy}, an archive, phonon labels.}
{\texttt{parity.npz}, \texttt{band\_parity.npz}, \texttt{irreps.npz}, figures.}
{Every state $\pm1$ (fractional values: a mixed near-degenerate pair; \texttt{--rotate} in \texttt{selection-rule});
forbidden/allowed at the level of numerical noise; every multiplet closes.}

\subsection{Optics}
\begin{cmd}
xphd dipoles --save <SAVE> --bse-dir <BSE_DIR> -o exc_dipoles.npy --deg-tol 0.005
xphd chirality --save <SAVE> --bse-dir <BSE_DIR> --dipoles exc_dipoles.npy --states 1 2
\end{cmd}
\step{Complex $x,y,z$ dipoles of the $\bQ=0$ states from one BSE (validated against yambo's residuals); helicity and valley
polarisation of a degenerate pair.}{\texttt{SAVE}, the BSE with \texttt{ndb.dipoles}.}{\texttt{exc\_dipoles.npy}; numbers.}
{Residuals agree per degenerate multiplet; a bright valley pair has helicity $\pm1$.}

\subsection{Transport}
\begin{cmd}
xphd bte '<ARCH>/GI_ExcPh_Q*.npz' --T 77 --hw-fine hw_fine.npy --parity parity.npz --window 0.1 --plan
xphd bte '<ARCH>/GI_ExcPh_Q*.npz' --T 77 --hw-fine hw_fine.npy --parity parity.npz --window 0.1 \
     --acoustic-model shell --acoustic-scaling omega --inject uniform --state 3 4 \
     --t-end 20000 --nt 200 --workers 8 --snapshots snaps_77K.npz
\end{cmd}
\step{Populations of every state at every $\bQ$ under the full rate matrix (detailed balance, out-rates refined on the fine mesh).
\texttt{--window} keeps the states within that energy of the minimum on the refined mesh (\texttt{--plan} prints them and the memory first).
\texttt{--inject}: \texttt{optical} (oscillator-weighted, optional \texttt{--pump-center/--pump-sigma}), \texttt{thermal}, or
\texttt{uniform} over \texttt{--state} (a degenerate bright pair: the basis-independent unpolarised pump).}
{Every archive of the zone, \texttt{hw\_fine.npy}; \texttt{parity.npz} with \texttt{--window}.}
{The snapshot file (populations, energies, per-state widths at every delay).}
{$N/N_0=1$ throughout; the final distribution against Boltzmann at $T$ (printed).}

\subsection{Photoemission}
\begin{cmd}
xphd arpes --save <SAVE> --bse-dir <BSE_DIR> --qp-db <GW>/ndb.QP --temperature 300 --nstates 100 -o arpes_300K.npz
xphd arpes --save <SAVE> --bse-dir <BSE_DIR> --qp-db <GW>/ndb.QP --snapshots snaps_77K.npz \
           --times 0 1000 10000 --path M K G "K'" M --sigma 0.01 --box K 0.12 --valence \
           -o arpes/arpes_{t}fs.npz --series arpes/series.npz
\end{cmd}
\step{$I(\mathbf{k},E)=\sum N_{S\bQ}\sum_{vc}|A^{S\bQ}_{vc\mathbf{k}}|^2\,\delta(E-E_{S\bQ}-\varepsilon_{v,\mathbf{k}-\bQ})$ with Boltzmann
(\texttt{--temperature}) or transport (\texttt{--snapshots}) occupations; maps on a path, momentum boxes, the valence band with its holes,
the spectra against the delay.}
{\texttt{SAVE}, the BSE, \texttt{ndb.QP}; a snapshot file.}
{One map per delay and the series file.}
{The replica lies below the conduction band by the binding energy; the zone-integrated intensity is constant in time.
With spin--orbit bands use \texttt{--spin-deg 1}. In PowerShell write \texttt{"K'"} in double quotes.}

% =====================================================================
\section{The run report}\label{sec:report}
After every command, also after a failure, rank 0 writes a report: \texttt{<output>\_report.pdf} next to the main
output, or, when the run wrote nothing, \texttt{xphd\_<command>\_<time>\_report.pdf} in the working directory (\texttt{.txt}
on Linux, below). The PDF holds:
\begin{itemize}[nosep,leftmargin=1.4em]
\item a title bar naming the step (BTE, ARPES, Linewidths, \ldots) with the version;
\item the xphd banner printed at the start of the run (version, developers and their affiliations, how to cite,
acknowledgements), beside a diagram of the pipeline (Quantum ESPRESSO, LetzElPhC, yambo and xphd) with this command's step marked;
\item the command line, status, start time and wall time, working directory;
\item the environment: xphd, Python and library versions, host, MPI ranks, thread settings;
\item every input file named on the command line with size, date and a SHA-256 checksum (files up to 64~MB);
\item every parameter, defaults included;
\item every file written or changed by the run;
\item plots of the main results, for the outputs of \texttt{generate}, \texttt{parity}, \texttt{linewidth}, \texttt{sweep},
\texttt{lineshape}, \texttt{bte} and \texttt{arpes}, each figure with lettered, titled panels and a caption saying what is
drawn; quantities over the zone are shown along $\Gamma$--M--K--$\Gamma$ (the mesh points on the path) and over the hexagonal first
Brillouin zone. Only the essential panels are drawn. An archive: the phonon dispersion, the final exciton bands with the one-phonon
window, the coupling out of the lowest state over the zone and per branch along the path. \texttt{parity}: the exciton bands coloured
by $\sigma_h$ parity. \texttt{linewidth}: the width of every state, and of state 1 by branch. \texttt{sweep}: the width of the lowest
exciton over the zone and along the path. \texttt{lineshape}: the width of the A line and $\varepsilon_2$. \texttt{bte}: the mean energy,
the energy distribution and the population over the zone. \texttt{arpes}: the map along the path and the occupied excitons
against their momentum; against the delay, the whole-zone spectrum and the mean exciton energy. Other files, such as
\texttt{Dmats.npy}, are listed and not drawn;
\item the full log as the terminal leaves it (a progress bar is one line, its final state), and the error if the run failed.
\end{itemize}
Off with \texttt{xphd --no-report <command> ...} or \texttt{XPHD\_REPORT=0}.

\textbf{On Linux and clusters.} The system decides the kind of report: the PDF on Windows and macOS, a text report,
\texttt{<output>\_report.txt}, on Linux (a cluster or an Ubuntu machine), where a terminal has no PDF viewer. The text
report holds the same banner, run record, inputs with checksums, parameters, outputs and full log, everything but the plots:
\begin{cmd}
less lw_report.txt                                  # on the cluster
\end{cmd}
\texttt{XPHD\_REPORT=pdf}, \texttt{txt} or \texttt{both} overrides the choice. With \texttt{both} on a cluster the PDF is drawn too
(no display is needed; DejaVu fonts where the desktop ones are missing) and the text report lists its figures by title, panels
and caption; copy it home for the plots (\texttt{scp user@cluster:/path/to/run/lw\_report.pdf .}, or WinSCP). The PDF needs
\texttt{matplotlib} (the \texttt{plot} extra, installed by \texttt{scripts/hpc/install\_cluster.sh}); without it the text report is written.

% =====================================================================
\section{Running in parallel}
The numbers do not depend on the layout: serial, process-pool and MPI runs write identical arrays.

{\small\renewcommand{\arraystretch}{1.2}
\begin{tabularx}{\linewidth}{>{\raggedright\arraybackslash}p{3.6cm}>{\raggedright\arraybackslash}p{5.4cm}>{\raggedright\arraybackslash}X}
\hdrow \hd{Command} & \hd{Spread over} & \hd{Memory per process}\\
\texttt{generate} & MPI ranks over $\bQ$; BLAS threads in each rank & one copy of the el--ph and BSE databases\\
\texttt{sweep} & MPI ranks (one per core) or \texttt{--workers} & one archive and the fine mesh\\
\texttt{bte}, out-rate refinement & MPI ranks or \texttt{--workers} & bounded by \texttt{--mem-per-task} / \texttt{XPHD\_TASK\_MEM\_GB}\\
\texttt{bte}, propagation & one process, every core as BLAS threads & about $10\times(\text{states})^2\times8$ bytes\\
every other command & rank 0 only & \\
\end{tabularx}}

\textbf{Install on a cluster} (\texttt{scripts/hpc/install\_cluster.sh} does this): \texttt{mpi4py} must be built against the
MPI that launches the job; the PyPI wheel brings its own MPI and runs unconnected copies.
\begin{cmd}
module load <compilers and MPI used for QE/yambo>
python3 -m venv ~/xphd-env && source ~/xphd-env/bin/activate
MPICC=$(which mpicc) pip install --no-binary=mpi4py mpi4py      # Intel MPI: MPICC=mpiicc
pip install -e "./xphd[hpc,yambo,plot]"
mpirun -np 4 xphd hpc-info                                      # ranks, nodes, cores, BLAS threads, memory
\end{cmd}
\textbf{Sizing.} \texttt{generate}: measure one $\bQ$ (\texttt{/usr/bin/time -v}, maximum resident size); ranks per node
$=\lfloor 0.85\times$ node memory / that size$\rfloor$, cores per rank as BLAS threads. \texttt{sweep} and the \texttt{bte} refinement:
one rank per core, \texttt{OMP\_NUM\_THREADS=1}; more ranks than $\bQ$ points gain nothing. \texttt{bte} propagation: one node,
all cores as threads (\texttt{--threads}); unbind it (\texttt{srun --cpu-bind=none}, \texttt{mpirun --bind-to none}).
\texttt{generate --resume} (default for several $\bQ$) and \texttt{--checkpoint DIR} (\texttt{sweep}, \texttt{bte}) make a stopped job
resumable: resubmit the same script. SLURM and PBS templates for every stage are in \texttt{scripts/hpc/}.

{\small\renewcommand{\arraystretch}{1.15}
\begin{tabularx}{\linewidth}{>{\ttfamily}p{5.0cm}X}
\hdrow \hd{Variable} & \hd{Meaning}\\
OMP\_NUM\_THREADS & BLAS threads per process (\texttt{--threads}, \texttt{--blas-threads} override)\\
XPHD\_TASK\_MEM\_GB & memory bound per \texttt{bte} refinement task (default 2)\\
XPHD\_MPI=0 / 1 & disable / force MPI detection\\
XPHD\_ALL\_RANKS=1 & every rank prints (debugging)\\
XPHD\_MP\_START & start method for \texttt{--workers}\\
XPHD\_NO\_BANNER=1 & no start-up banner\\
XPHD\_REPORT=0 & no report\\
XPHD\_REPORT=pdf / txt / both & which report; by default the PDF on Windows and macOS, the text report on Linux\\
\end{tabularx}}

% =====================================================================
\section{When a result looks wrong}
{\small\renewcommand{\arraystretch}{1.2}
\begin{longtable}{>{\raggedright\arraybackslash}p{4.6cm}>{\raggedright\arraybackslash}p{5.0cm}>{\raggedright\arraybackslash}p{6.6cm}}
\hdrow \hd{Symptom} & \hd{Cause} & \hd{Fix}\\
\endhead
Every BSE energy exactly 0, eigenvectors fine & a compiler fault in yambo's \texttt{BS\_E=cmplx(BS\_E\_real,...)} (Intel ifort) & the element loop in \texttt{K\_diago\_driver.F}\\
LetzElPhC: ``Unable to open the dyn0 file'' & \texttt{ph\_save\_dir} not set & \texttt{ph\_save\_dir = ph.save}\\
dvscf files of 0 bytes after a \texttt{ph.x} restart & \texttt{recover} trusted old \texttt{phsave} records & delete those \texttt{dynmat.<iq>.*.xml}, rerun the $\bq$ without \texttt{recover}\\
\texttt{check-archive}: one of TR/$C_3$ far above the other & a rotation fault, or a doublet split by long-range exchange & \texttt{--detail}; raise \texttt{--deg-tol}\\
Fractional parities in a ring of $\bQ$ & mixed near-degenerate states of opposite parity & \texttt{parity --deg-tol}, \texttt{selection-rule --rotate}\\
Width at a band minimum far too small & the acoustic channel inside the first cell is missing & \texttt{--acoustic-model q2} (\texttt{linewidth}) or \texttt{shell} (\texttt{lineshape}, \texttt{bte})\\
Band widths too steep near crossings & global Fourier interpolation rings & \texttt{--band-interp local} (default)\\
\texttt{matdyn} outputs broken after \texttt{cat} & a header and cumulative distances per chunk & \texttt{xphd matdyn-merge}\\
$\varepsilon_2$ line much narrower than its width at $E_A$ & one-shot $\Sigma$ at the band bottom & \texttt{lineshape --final-width sc} (default)\\
yambopy ``open-shell'' error reading \texttt{ndb.dipoles} & band count when the first band is not 1 & \texttt{xphd dipoles}\\
\texttt{hpc-info}: ``MPI not active'' under \texttt{srun} & \texttt{mpi4py} not built against this MPI, or \texttt{srun} without \texttt{--mpi=pmix} & rebuild \texttt{mpi4py} from source; \texttt{srun --mpi=pmix}\\
``launched with N MPI ranks but mpi4py cannot be imported'' & the environment is not activated in the job script & \texttt{source \textasciitilde/xphd-env/bin/activate} in the script\\
K/K$'$ weights differ transiently in \texttt{bte} & the solver's basis inside a degenerate pair & \texttt{--inject uniform --state <pair>}, \texttt{--deg-tol}\\
\end{longtable}}

\iffull
\input{extra}
\fi

\appendix
\section{Command reference}
Generated from \texttt{xphd <command> --help}.
\input{reference}

\end{document}
